% Sablon pentru realizarea lucrarii de licenta, conform cu recomandarile
% din ghidul de redactare:
% - https://fmi.unibuc.ro/finalizare-studii/
% - https://drive.google.com/file/d/1xj9kZZgTkcKMJkMLRuoYRgLQ1O8CX0mv/view

% Multumiri lui Gabriel Majeri, acest sablon a fost creat pe baza
% codului sursa a lucrarii sale de licenta. 
% Codul sursa: https://github.com/GabrielMajeri/bachelors-thesis
% Website: https://www.gabrielmajeri.ro/
%
% Aceast sablon este licentiat sub Creative Commons Attribution 4.0 International License.

\documentclass[12pt, a4paper]{report}

% Suport pentru diacritice și alte simboluri
\usepackage{fontspec}

\usepackage{lipsum}

% Suport pentru mai multe limbi
\usepackage{polyglossia}

% Setează limba textului la română
\setdefaultlanguage{romanian}
% Am nevoie de engleză pentru rezumat
\setotherlanguages{english}

% Indentează și primul paragraf al fiecărei noi secțiuni
\SetLanguageKeys{romanian}{indentfirst=true}

% Suport pentru diferite stiluri de ghilimele
\usepackage{csquotes}

\DeclareQuoteStyle{romanian}
  {\quotedblbase}
  {\textquotedblright}
  {\guillemotleft}
  {\guillemotright}

% Utilizează biblatex pentru referințe bibliografice
\usepackage[
    maxbibnames=50,
    sorting=nty
]{biblatex}

\addbibresource{bibliography.bib}

% Setează spațiere inter-linie la 1.5
\usepackage{setspace}
\onehalfspacing

% Modificarea geometriei paginii
\usepackage{geometry}

% Include funcțiile de grafică
\usepackage{graphicx}
% Încarcă imaginile din directorul `images`
\graphicspath{{./images/}}

% Listări de cod
\usepackage{listings}

% Linkuri interactive în PDF
\usepackage[
    colorlinks,
    linkcolor={black},
    menucolor={black},
    citecolor={black},
    urlcolor={blue}
]{hyperref}

% Comenzi matematice
\usepackage{amsmath}
\usepackage{mathtools}

% Simboluri matematice codificate Unicode
\usepackage[warnings-off={mathtools-colon,mathtools-overbracket}]{unicode-math}

% Formule matematice
\newcommand{\bigO}[1]{\symcal{O}\left(#1\right)}
\DeclarePairedDelimiter\abs{\lvert}{\rvert}

% Lista de ecuatii
\usepackage{tocloft}

\newcommand{\listequationsname}{Listă de ecuații}
\newlistof{myequations}{equ}{\listequationsname}
\newcommand{\myequations}[1]{%
\addcontentsline{equ}{myequations}{\protect\numberline{\theequation}#1}\par}

% Suport pentru rezumat în două limbi
% Bazat pe https://tex.stackexchange.com/a/70818
\newenvironment{abstractpage}
  {\cleardoublepage\thispagestyle{empty}}
  {\vfill\cleardoublepage}
\renewenvironment{abstract}[1]
  {\bigskip\selectlanguage{#1}%
   \begin{center}\bfseries\abstractname\end{center}}
  {\par\bigskip}

% Suport pentru anexe
\usepackage{appendix}

% Stiluri diferite de headere și footere
\usepackage{fancyhdr}

% Metadate
\title{Emotion-Aware Desktop AI Assistant with Facial Emotion Recognition}
\author{Preda Alexandru-Florin}

% Generează variabilele cu @
\makeatletter

\begin{document}

% Front matter
\cleardoublepage
\let\ps@plain

% Pagina de titlu
\include{0-title}
\restoregeometry
\newgeometry{
    margin=2.5cm
}

\fancypagestyle{main}{
  \fancyhf{}
  \renewcommand\headrulewidth{0pt}
  \fancyhead[C]{}
  \fancyfoot[C]{\thepage}
}

\addtocounter{page}{1}

% Rezumatul
\begin{abstractpage}

\begin{abstract}{romanian}


Scopul acestei lucrări este de a prezenta proiectarea și implementarea unui \textbf{asistent AI desktop conștient de emoții, cu recunoaștere facială a emoțiilor}. Aplicația îi permite utilizatorului să interacționeze cu asistentul prin text și voce, folosind în același timp recunoașterea facială a emoțiilor pentru a îmbunătăți contextul interacțiunii. În loc să se bazeze doar pe ceea ce utilizatorul scrie sau spune, asistentul ia în considerare și modul în care utilizatorul pare să se simtă în timpul conversației.

Pentru a atinge acest obiectiv, aplicația este construită ca un sistem modular care conectează mai multe componente de inteligență artificială într-un singur asistent desktop. Sistemul combină procesarea intrărilor utilizatorului, analiza emoțiilor faciale, transcrierea vorbirii, generarea răspunsurilor și comunicarea în timp real între module. Starea emoțională detectată este folosită ca un context suplimentar pentru asistent, permițând ca răspunsurile generate să fie mai bine adaptate la interacțiunea curentă.

Soluția este formată din cinci componente principale:

\begin{itemize}
    \item O interfață de asistent desktop care îi permite utilizatorului să interacționeze cu sistemul prin text și voce;

    \item O componentă backend responsabilă de gestionarea sesiunii de interacțiune și a comunicării dintre modulele aplicației;

    \item O componentă speech-to-text utilizată pentru convertirea intrării vocale a utilizatorului în text;

    \item O componentă de recunoaștere facială a emoțiilor utilizată pentru detectarea stării emoționale aparente a utilizatorului din imaginile capturate de webcam;

    \item O componentă de generare a răspunsurilor care primește mesajul utilizatorului împreună cu contextul emoțional detectat și generează răspunsul asistentului.
\end{itemize}

Lucrarea se concentrează pe prezentarea celor cinci componente și a modului în care acestea sunt conectate în aplicația finală, precum și pe antrenarea și evaluarea modelului de recunoaștere facială a emoțiilor, integrarea contextului emoțional în răspunsurile asistentului și prezentarea rezultatelor experimentale obținute.

\end{abstract}

\newpage

\begin{abstract}{english}

The goal of this paper is to present the design and implementation of an \textbf{Emotion-Aware Desktop AI Assistant with Facial Emotion Recognition}. The application allows the user to interact with an assistant through text and voice, while also using facial emotion recognition to improve the context of the interaction. Instead of relying only on what the user writes or says, the assistant also takes into account how the user appears to feel during the conversation.

To achieve this goal, the application is built as a modular system that connects several artificial intelligence components into a single desktop assistant. The system combines user input processing, facial emotion analysis, speech transcription, response generation and real-time communication between the modules. The detected emotional state is used as additional context for the assistant, allowing the generated answers to be more adapted to the current interaction.

The solution consists of five main components:

\begin{itemize}
    \item A desktop assistant interface that allows the user to interact with the system through text and voice;

    \item A backend component responsible for managing the interaction session and the communication between the application modules;

    \item A speech-to-text component used for converting the user's voice input into text;

    \item A facial emotion recognition component used for detecting the user's apparent emotional state from webcam images;

    \item A response generation component that receives the user's message together with the detected emotional context and generates the assistant's answer.
\end{itemize}

The paper focuses on the presentation of the five components and their connection in the final application, as well as the training and evaluation of the facial emotion recognition model, the integration of emotional context into the assistant's responses, and the presentation of the experimental results obtained.

\end{abstract}

\end{abstractpage}


% Cuprins
\tableofcontents
\newpage
\listoffigures
\newpage
\listoftables
\newpage
\listofmyequations

% Main matter
\cleardoublepage
\pagestyle{main}
\let\ps@plain\ps@main

\chapter{Introduction}

\section{Motivation}

In the current era of AI assistants, most human-computer interactions are based on text input, voice input, or text-to-speech. When people communicate with each other, they also use non-verbal signals, such as facial expressions, tone of voice, hesitation, and body language. A desktop assistant that can observe the apparent facial emotion can use this extra context to respond in a more natural way.

These interactions often contain sensitive information, such as personal messages, voice recordings, and webcam images. For this reason, running the main components locally is valuable, because it gives the user more control over their data.

In my project, I explore whether facial emotion recognition can be useful in the interaction between humans and AI assistants, not only as a standalone image classification task.

\section{Problem Statement}

Standard interaction between users and AI assistants is usually based only on explicit input, such as written or spoken text. Because of this, the assistant is not aware of the user's emotional state.

Although current classification models have evolved from classical CNN architectures to deeper models and architectures that may include attention mechanisms, facial emotion recognition remains a difficult task. In real-world webcam conditions, predictions can be affected by lighting, camera quality, face position, occlusions, and the fact that facial expressions do not always reflect the user's real emotion.

For this reason, the detected facial emotion should not be treated as an absolute truth. A separate challenge is to use this prediction only as additional context in the interaction between the user and the assistant.

The main challenge of my project is to integrate several AI components into one local architecture that can run on a regular user's device: webcam-based emotion recognition, speech-to-text, LLM response generation, session management, and real-time response streaming.


\section{Objectives and Personal Contributions}

The main goal of this dissertation is to design, implement, and evaluate an emotion-aware desktop AI assistant that supports both text and voice interaction. The system uses the user's apparent facial emotion as additional context for a locally running LLM.

To achieve this, I developed a modular architecture based on a FastAPI backend. The backend exposes separate endpoints for session management, emotion prediction, speech transcription, and assistant response generation. This separation makes the system easier to test, extend, and connect to different user interfaces.

The main contribution of the project is the facial emotion recognition component. I implemented a PyTorch training pipeline based on MobileNetV2 and trained the model on a public facial expression dataset. The pipeline contains the full training process, from preparing the images and handling class imbalance, to saving the trained model and evaluating its performance.

The experimental part of the dissertation focuses on finding the best training setup for the emotion recognition model. I started by testing how class imbalance should be handled, then checked how image quality filtering affects the results. After that, I used the findings to tune the model parameters, compare baseline and fine-tuned versions, and train the final model on the full dataset.

Another contribution is the integration of the trained emotion model into the assistant workflow. The webcam feed is processed periodically, the predicted emotion is stored in the current session, and this emotional context is inserted into the prompt sent to the local LLM. The system treats the predicted emotion as a soft signal, not as a certain representation of the user's real emotional state.

The application also includes voice interaction through a speech-to-text model. Microphone input is recorded, transcribed, and sent to the assistant as text. The final response is generated through Ollama and streamed back to the user in real time.

Overall, the project combines machine learning experimentation with the implementation of a working end-to-end assistant system.

\section{Structure of the Thesis}

The thesis starts with an introduction to the project and to the motivation behind it. This first chapter explains why emotion-aware interaction can be useful for AI assistants, what problem the dissertation addresses, and what I wanted to achieve through the implementation.

The second chapter presents the background needed to understand the project. It introduces affective computing and facial emotion recognition, then discusses public datasets used for this task. It also presents existing tools that can be used for building local AI assistants, such as speech-to-text models, local language models, and computer vision libraries.

The third chapter focuses on the facial emotion recognition model. It describes the dataset used for training, the preprocessing steps, the MobileNetV2 architecture, and the training pipeline. This chapter also presents the experiments used for selecting the final model, including the comparison of imbalance handling methods, quality filtering, baseline training, and fine-tuning.

The fourth chapter presents the system design and architecture of the final application. It explains how the desktop interface, FastAPI backend, session management, emotion recognition service, speech-to-text service, prompt construction, and local LLM are connected in the full assistant workflow.

The final chapter summarizes the work done in the dissertation and presents the main conclusions. It also discusses the limitations of the current implementation and possible directions for improving the system in the future.

\chapter{Background and Related Work}

\section{Affective Computing and Facial Emotion Recognition}

Affective computing is a research area focused on systems that can recognize, interpret, and react to human emotions. The field was introduced by Picard, who argued that emotion is an important part of intelligent human-computer interaction \cite{picard1997affective}.

In real communication, people do not use only words. They also use facial expressions, voice tone, gestures, and context. Because of this, emotion-aware systems may analyze several types of input, such as face images, speech, and text \cite{poria2017review,mobbs2025emotionreview}. This dissertation focuses mainly on facial emotion recognition, while speech is used only as a way of converting the user's voice into text.

Facial Emotion Recognition is the task of predicting a person's apparent emotion from an image or video frame of the face. Common emotion classes include anger, disgust, fear, happiness, sadness, surprise, and neutrality. In my project, the detected facial emotion is not treated as a certain description of the user's real emotional state. Instead, it is used as additional context for the assistant.

This distinction is important because facial emotion recognition is a difficult task. In practical webcam conditions, predictions can be affected by lighting, camera quality, face position, occlusions, and the ambiguity of facial expressions. Recent surveys also show that performance depends strongly on the dataset, the model architecture, and the evaluation setting \cite{mobbs2025emotionreview}. For this reason, the assistant uses the predicted emotion as a soft signal, not as a rule that controls the direction of generated answer.

\section{Public Datasets and Benchmarks for Facial Emotion Recognition}

Facial emotion recognition models are usually trained and compared on public datasets. One of the most used benchmarks is FER2013, introduced in the Challenges in Representation Learning competition \cite{goodfellow2015challenges}. FER2013 contains low-resolution facial images labeled with seven emotion classes: angry, disgust, fear, happy, sad, surprise, and neutral.

FER2013 is useful because it is public and widely used, which makes comparison with other work easier. At the same time, it is not an easy dataset. The images have low resolution, the labels can be noisy, and the classes are not perfectly balanced. These problems are important because they influence both training and evaluation.

Other datasets also exist for facial emotion recognition. AffectNet is a larger dataset that contains both categorical emotion labels and valence-arousal annotations \cite{mollahosseini2019affectnet}. Other examples include CK+, JAFFE, RAF-DB, and FERPlus. These datasets differ in size, image quality, annotation method, and acquisition conditions.

Recent work still uses FER2013 to evaluate modern architectures. For example, Naik et al. use an EfficientNetB2-based model and report results on FER2013, while also pointing out problems such as low image quality, label noise, and class imbalance \cite{naik2026efficientnetb2}. This makes FER2013 a relevant choice for this dissertation, especially because the goal is not only classification accuracy, but also integration into a real-time assistant.

\section{Existing Components for Local AI Assistants}

Open-source tools make it possible to run AI assistant components directly on a personal computer. This is important when the application works with private data, such as user messages, microphone recordings, or images from a webcam.

For speech-to-text, one of the most widely used open-source options is Whisper, a pretrained model for automatic speech recognition \cite{radford2022whisper}. Other available alternatives are wav2vec 2.0, which is based on self-supervised speech representation learning, and Vosk, which provides offline speech recognition models \cite{baevski2020wav2vec}.

For response generation, large language model families provide versions that can be executed locally. Examples include Llama, Gemma, Mistral, Qwen, and Phi \cite{touvron2023llama2,gemma2024,jiang2023mistral}. Tools such as Ollama simplify local deployment by downloading models, exposing them through a local API, and handling inference on the user's machine.

There are also local libraries and frameworks for computer vision tasks. PyTorch and torchvision are commonly used for training and running neural networks, while OpenCV is often used for webcam access, face detection, frame processing, and image manipulation. These tools provide the basic building blocks required for applications that combine speech, language, and vision on the same device.

\chapter{System Design and Architecture}

\section{Functional and Non-Functional Requirements}

\section{High-Level Architecture}

\section{Component Overview}

\subsection{Desktop Assistant Interface}

\subsection{Backend API Layer}

\subsection{Emotion Recognition Service}

\subsection{Speech-to-Text Service}

\subsection{LLM Generation Service}

\section{Data Flow During a User Interaction}

\section{Session and Emotional State Management}

\section{API Design and Service Boundaries}

\section{Privacy and Local Execution Considerations}

\chapter{Facial Emotion Recognition Model}

\section{Dataset Description}

\section{Data Preprocessing and Augmentation}

\section{Class Imbalance in FER2013}

\section{Model Architecture}

\section{Training Strategy}

\subsection{Baseline Training}

\subsection{Fine-Tuning}

\section{Evaluation Metrics}

\section{Checkpointing and Experiment Tracking}

\chapter{Experimental Methodology}

\section{Experimental Goals}

\section{Hardware and Software Environment}

\section{Experiment 1: Comparing Imbalance Tactics}

\section{Experiment 2: Comparing Quality Filtering}

\section{Experiment 3: Baseline Sweep with Injected Settings}

\section{Experiment 4: Fine-Tuning Sweep with Injected Settings}

\section{Model Selection Procedure}

\chapter{Results and Discussion}

\section{Imbalance Strategy Comparison}

\section{Quality Filtering Comparison}

\section{Baseline Sweep Results}

\section{Fine-Tuning Sweep Results}

\section{Best Model Analysis}

\section{Confusion Matrix and Per-Class Performance}

\section{Error Analysis}

\section{Limitations}

\chapter{Assistant Implementation and Integration}

\section{FastAPI Backend Implementation}

\section{Emotion Prediction Endpoint}

\section{Real-Time Webcam Emotion Polling}

\section{Speech-to-Text Endpoint}

\section{LLM Response Generation with Ollama}

\section{Prompt Construction with Emotional Context}

\section{Streaming Responses with Server-Sent Events}

\section{End-to-End Interaction Flow}

\chapter{Testing and Validation}

\section{Backend Endpoint Testing}

\section{Session Management Tests}

\section{Assistant Response Streaming Tests}

\section{Manual Validation with Webcam and Voice Input}

\section{Demo Scenarios}

\chapter{Conclusions and Future Work}

\section{Summary of the Work}

\section{Main Contributions}

\section{Conclusions from the Experiments}

\section{Limitations}

\section{Future Improvements}


\appendix
\chapter{API Endpoint Reference}

\section{Session Endpoints}

\section{Emotion Recognition Endpoints}

\section{Speech-to-Text Endpoints}

\section{Assistant Response Endpoint}

\chapter{Hyperparameter Sweep Configurations}

\section{Imbalance Strategy Runs}

\section{Quality Filtering Runs}

\section{Baseline Sweep Configuration}

\section{Fine-Tuning Sweep Configuration}

\chapter{Additional Result Tables}

\section{Training Runs}

\section{Validation Metrics}

\section{Test Metrics}

\chapter{Additional Confusion Matrices}

\section{Baseline Models}

\section{Fine-Tuned Models}

\chapter{Installation and Running Instructions}

\section{Backend Setup}

\section{Model Checkpoint Configuration}

\section{Running the Desktop Assistant}

\section{Running the Experiment Scripts}

\chapter{Selected Code Listings}

\section{Backend Application Factory}

\section{Emotion Prediction Service}

\section{Training Loop}

\section{Prompt Construction}


\printbibliography[heading=bibintoc]

\end{document}
